\subsection{平展映射的局部截面}

设 E 与 B 为集合，A 为 B 的一个子集，\(p: E \rightarrow B\) 为一个映射。p 在 A 上的一个截面是一个映射 \(s: A \rightarrow E\) ，使得 \(p \circ s\) 是 A 到 B 中的典型单射。给出 p 在 A 上的一个截面 s，等价于给出由 p 诱导的映射 \(p_{\mathrm{A}} \colon \overline{p}^{1}(\mathrm{A}) \to \mathrm{A}\) 的一个截面。若 s 是 p 在 A 上的一个截面且 \(A'\) 是 A 的一个子集，则限制 \(s'\) （s 到 \(A'\) 上的）是 p 在 \(A'\) 上的一个截面。这时我们称 s 是 \(s'\) 到 A 上的一个延拓。

当 E 与 B 是拓扑空间且 \(p \colon \mathrm{E} \to \mathrm{B}\) 是连续映射时，集合 \(\mathcal{C}_{\mathrm{B}}(\mathrm{A};\mathrm{E})\) (I, p. 2)（即 \(p\) 在 A 上的连续截面所成的集合）也记作 \(\mathcal{S}(\mathrm{A};p)\) 或 \(\mathcal{S}(\mathrm{A};\mathrm{E})\) 。设 \(s\) 为 \(p\) 在 A 上的一个连续截面。映射 \(s\) 诱导出从 A 到 \(s(\mathrm{A})\) 上的一个同胚，而 \(p\) 诱导出其逆同胚。于是由 I, p. 28 的定义 2，我们有：

命题 9. --- 设 p: E → B 为一个连续映射。为使 p 是平展的，必须且只需 E 的每个点都有一个开邻域，它是在 B 的某个开子集上定义的 p 的一个连续截面的像。

注. --- 1) 设 \(p \colon \mathrm{E} \to \mathrm{B}\) 为一个连续且开的映射。为使 E 的一个开集 U 是 \(p\) 在 B 的某个开集上的一个连续截面的像，必须且只需 \(p\) 在 U 上的限制是单射的。

\begin{enumerate}
\def\labelenumi{\arabic{enumi})}
\setcounter{enumi}{1}
\tightlist
\item
  设 \(p \colon \mathrm{E} \to \mathrm{B}\) 为一个连续映射。设 B 的每个点都有一个开邻域 V，使得存在 \(p\) 在 V 上的一个连续截面。这样的映射 \(p\) 未必是平展的，甚至未必是开的。然而它是满射且普遍严格的 (I, p. 20, 推论)。
\end{enumerate}

